The model must be well-defined to be used in NMPC, in such a way 
to obtain reliable predictions of system behavior in response to various control actions.
The usage of a data-driven approach, such as neural networks, in the vehicle modeling phase has the
advantage of not modeling all the nonlinear physical effects, which can be also very complex.
In fact, neural networks have the capability to accurately model them. 
Neural networks are universal approximators, and they can be used both for designing 
the entire dynamic model of a system and approximating the disturbances and noises 
a system can be affected.
Our purpose is to approximate the entire kinematic model of a quadrotor with a neural model, by first testing the performances on a unicycle model due to its low complexity. 
In this section we will start describing briefly the kinematic model of a unicycle, and then 
we will describe the neural architecture used to approximate it.

\subsubsection{Architecture}
The general architecture of the networks is a MLP with as input the state of the system, the input and (in the PINN models) the time stamp.
We tested the proposed approach with a various number of hidden layers and neurons, keeping the dimension of the network reasonably small 
for enanche the possibility the on-board computer of the Quadrotor.
\\

In the standard neural network, a smaller number of neurons was sufficient for exploit the dynamics, while the PINN framework required a larger number of neurons for smooth 
the loss function. Having a smoother loss landscape is crucial for let the model mimic the real dynamics.

In general, the combination of small hidden state and the absence of regularization terms in the standard models are 
the reason for which the general nerual network framework is not able to approximate the underlying physics laws over the entire training workspace.


\subsection{Standard neural network}

\subsubsection{Training scheme}

As can be noticed in the previous section, the dataset does not contain the dynamics of the system, 
but just the evolution of the states. On the contrary, the forward pass of the network returns
the vehicle dynamics. So, what we did is to pass the first pair $(s_t, a_t)$ of a sample to the 
network; then we integrated the output in the sample time $dt$, 
and added to the previous state (which is also the input of the network).
In this way we obtain the state $s_{t+1}$.
At this point we concatenate the new obtained state with the action $a_{t+1}$ given by the sample dataset
and we pass it to the network again, and repeating
this process $k$ times we get the states evolution according to the neural model. \\
Ended this mechanism, the loss is computed as the mean-squared error between the neural and the nominal evolutions. 

The multi-step training is a crucial step for a non-physics informed model to learn the total evolution. 
Experiment in a single step prediction has shown that the displacement of the vehicle is to small for having a fine-grain error measurement. 

\subsubsection{Consideration of the Dataset}
As anticiapted before, the dataset for such approach do not require any particular constraint. This allow us to test standard normalization techniques
and so robustify the model to different input scales. 
For doing so, we used a min-max normalization over the input, but we can also incorporate over the MLP architecture a normalization layer, that increase the stability of the train.

We prefered to use the layer norm in our experiment because such module do not require a preprocessing at inference time, before invoke the MPC framework.
However, the normalization layer is still trained over the dataset, limiting the capacity of the network over range outside the training one. 

\subsubsection{Unicycle}
In the application of the first approach over the unicycle model, the network was trained with 2 hidden state and a latent dimension of 128. 
The input was the concatenation of the state $x \in R^3$ and the action $u \in R^2$, while the output was the state derivative $\dot x \in R^3$.

\subsubsection{Quadrotor}
The quadrotor dynamic present a more complex model, with a state $x \in R^9$ and an action $u \in R^4$. For approximate such dynamic with the defined training scheme, 
a larger network as been tested: we increase the number of hidden layers up to 4. We have also tested the network with a latent dimension of 256, but the results were not competitive. 
In particular, considering the training time and the computational power required for the inference, we prefered to keep the same latent dimension of the unicycle network.


\subsection{Unicycle Neural model}
\todo[inline]{Aggiungere figura MLP}
\subsubsection{Architecture}
The neural network used to model the kinematic model of the unicycle is a simple MLP architecture,
with two hidden layers of 128 neurons. The network has two input which are concatenated together,
the current state of the unicycle (dimension 3) and the current applied actions on the vehicle (dimension 2),
so it must be of dimension 5 and is represented by 
\begin{equation}
    X = [x, y, \theta, v, \omega]
\end{equation}

The output of the network has dimension 3, and it represents the dynamics of the system, therefore it reproduces 
the differential equations which describe the motion of the unicycle:

\begin{equation}
    Y = [\dot{x}, \dot{y}, \dot{\theta}]
\end{equation}

\subsubsection{Dataset}
The dataset was generated by using the nominal model of the robot and generating 500 trajectories.
Each trajectory is composed by 1000 steps, where the inputs are maintained constant for the whole episode.
This choice was taken because it was noticed that the interval between
consecutive time-steps is too small, therefore the changes about the state of the unicycle are small too, and 
the network was not able to approximate correctly the motion of the vehicle. On the contrary, having the same
input for a lot of steps allows to understand in a better way the dynamics of the system subject to specific 
inputs. \\
Each sample is not just a pair state-action at step $t$, but it is a window of dimension $k$ which contains
all pairs state-action from time-step $t$ till $t+k$ (less the last action $a_{t+k}$ which is useless), 
so it is in the form 
\begin{equation}
    (s_t, a_t, s_{t+1}, a_{t+1}, ..., s_{t+k-1}, a_{t+k-1}, s_{t+k})
\end{equation}
This allows to have a larger window of evolution of the system with respect to just the motion executed in
one single step. 


\subsubsection{Training}

As can be noticed in the previous section, the dataset does not contain the dynamics of the system, 
but just the evolution of the states. On the contrary, the forward pass of the network returns
the vehicle dynamics. So, what we did is to pass the first pair $(s_t, a_t)$ of a sample to the 
network; then we integrated the output in the sample time $dt$, 
and added to the previous state (which is also the input of the network).
In this way we obtain the state $s_{t+1}$.
At this point we concatenate the new obtained state with the action $a_{t+1}$ given by the sample dataset
and we pass it to the network again, and repeating
this process $k$ times we get the states evolution according to the neural model. \\
Ended this mechanism, the loss is computed as the mean-squared error between the neural and the nominal evolutions. 


\subsection{Quadrotor Neural model}